\section{Conclusion}
\label{sec:conclusion}

This paper has established three findings about compactness in algorithmic redistricting at state legislative scale.

First, the mathematical relationship between district count and compactness is confirmed empirically: state house maps, which divide states into 40--400 districts, achieve systematically higher Polsby-Popper ratios than congressional maps, which divide states into 1--53 districts. The mean advantage is 0.027 PP units (7.5\%) at tract resolution, consistent with the $O(1/\sqrt{k})$ scaling law derived from the geometry of geographic partitioning.

Second, geographic resolution matters: moving from census tract to block group resolution improves mean state house PP by an additional 0.013 PP units (3.4\%). The resolution effect is largest in states with irregular geographic features (coastlines, river systems, mountain ridges) that block groups can trace more precisely than tracts. The combined scale and resolution effect produces state house PP at block group resolution (mean 0.401) that is 11\% higher than congressional PP at tract resolution (mean 0.361).

Third, algorithmic state house maps outperform enacted state house maps by a mean of 0.039 PP units (17\% improvement, significant at $p < 0.001$) across 35 states with available enacted map data. The improvement is largest in states where courts have identified gerrymandering: Maryland (+0.133), North Carolina (+0.131), Pennsylvania (+0.134). In these states, the enacted maps' low compactness reflects the non-compact district shapes required to achieve partisan objectives; the algorithmic maps, which cannot access partisan data, produce substantially more compact configurations.

The practical implications are threefold. For state legislatures considering algorithmic adoption, the compactness evidence is strongest at state legislative scale: the algorithm produces maps that are not merely competitive with enacted maps but substantially better on the criterion that most redistricting laws explicitly mandate. For courts evaluating redistricting challenges, the compactness comparison provides a straightforward standard: maps that fall substantially below the algorithmic benchmark (particularly below enacted PP minus one standard deviation) are candidates for scrutiny. For algorithmic designers, the finding that the algorithm maintains its compactness advantage at all chamber sizes confirms that METIS recursive bisection is a robust approach at the scales required for state legislative redistricting.

The principal limitation of this analysis is that compactness is not the only redistricting criterion. A state house map that is highly compact may still fail to preserve communities of interest, may split municipalities unnecessarily, or may fail VRA requirements. The compactness advantage established here is one dimension of a multi-dimensional evaluation. Future work should evaluate algorithmic state house maps on the full range of state-specific criteria, using the parameterization framework developed in the companion paper on state criteria variation.\footnote{See \citet{dellaLibera2026legal} (F.4).}

A second limitation is that the enacted map comparison is based on maps drawn under the pre-reform processes in many states; several of the states with the largest algorithmic advantage have since reformed their redistricting processes through commissions or court intervention. Evaluating algorithmic maps against commission-drawn maps, rather than legislative gerrymanders, would be a more demanding comparison and is left for future work.
